\documentclass[11pt]{article}
\usepackage{amsmath,amssymb,amsthm}
\usepackage[margin=2.8cm]{geometry}
\emergencystretch=3em

\newtheorem{theorem}{Theorem}
\newtheorem{proposition}[theorem]{Proposition}
\newtheorem{lemma}[theorem]{Lemma}
\newtheorem{corollary}[theorem]{Corollary}
\theoremstyle{definition}
\newtheorem{definition}[theorem]{Definition}
\newtheorem{example}[theorem]{Example}
\theoremstyle{remark}
\newtheorem{remark}[theorem]{Remark}

\newcommand{\abs}[1]{\lvert #1\rvert}
\newcommand{\norm}[2]{\lVert #1\rVert_{#2}}
\newcommand{\cJ}{\mathcal J}
\newcommand{\Z}{\mathbb Z}
\newcommand{\R}{\mathbb R}
\newcommand{\one}{\mathbf 1}
\newcommand{\upd}[3]{#1[#2 := #3]}

\title{Janson's Lemma~2, direction (ii)\,\(\Rightarrow\)\,(i)\\
\large (a proof plan for \texttt{mem\_\(\Omega\)\_of\_enorm\_le\_prod})}
\date{}

\begin{document}
\maketitle

\section{Setting and statement}

\paragraph{Quasinorms and moduli.}
A quasinorm \(X\) on a commutative monoid \(\gamma\) is a map
\(\norm{\cdot}{X}\colon \gamma \to [0,\infty]\) with \(\norm{0}{X} = 0\) and
\(\norm{x+y}{X} \le C_X(\norm{x}{X} + \norm{y}{X})\) for a constant \(1 \le C_X < \infty\)
(\texttt{EQuasinorm}, \(C_X\) = \texttt{X.C}). A modulus is a map
\(\abs{\cdot}\colon\gamma\to M\) into a commutative monoid \(M\) with a preorder. It is
\emph{decomposable} (\texttt{Modulus.IsDecomposable}) if \(\abs{0} = 0 \le \abs{x}\),
\(\abs{x+y} \le \abs{x} + \abs{y}\), and \(\abs{x} \le \abs{y_0} + \abs{y_1}\) implies
\(x = x_0 + x_1\) with \(\abs{x_j} \le \abs{y_j}\) and \(\abs{x_j} \le \abs{x}\). A quasinorm
\(X\) is \emph{solid} (\texttt{IsSolid}) if \(\abs{x} \le \abs{y}\) implies
\(\norm{x}{X} \le \norm{y}{X}\). Only transitivity of \(\le\) on \(M\) is available; in
particular two inequalities in \(M\) cannot be added.

\paragraph{The \(K\)- and \(J\)-functionals.}
Let \(X = (X_0, X_1)\) be a couple of quasinorms on \(\gamma\) (\texttt{Couple}). For
\(t \in [0,\infty]\) and \(x \in \gamma\),
\[
  K(t,x) = \inf_{x = x_0 + x_1}\bigl(\norm{x_0}{X_0} + t\norm{x_1}{X_1}\bigr)
  \quad(\texttt{kNorm}),
  \qquad
  J(t,x) = \max\bigl(\norm{x}{X_0},\, t\norm{x}{X_1}\bigr)
  \quad(\texttt{jNorm}),
\]
and \(\Delta(X) = \{x : \norm{x}{X_0} < \infty,\ \norm{x}{X_1} < \infty\}\), which is
\texttt{X.inf.finiteLocus} by \texttt{mem\_finiteLocus\_inf}. For \(\theta \in \R\),
\(q \in [0,\infty]\), \(\varphi\colon[0,\infty]\to[0,\infty]\) and
\(f\colon\Z\to[0,\infty]\) put
\begin{align*}
  \Phi_{\theta,q}(\varphi) &= \bigl\lVert t \mapsto t^{-\theta}\varphi(t)\bigr\rVert_{L^q((0,\infty),\,dt/t)}
  && (\texttt{phiFunctional}),\\
  \Phi^{\Z}_{\theta,q}(f) &= \bigl\lVert k \mapsto 2^{-k\theta} f_k\bigr\rVert_{\ell^q(\Z)}
  && (\texttt{discretePhiFunctional}).
\end{align*}
Both are defined through \texttt{eLpNorm}. The interpolation norms are
\begin{align*}
  \norm{x}{X_{\theta,q}} &= \Phi_{\theta,q}\bigl(K(\cdot,x)\bigr)
  && (\texttt{X.kMethod}\ \theta\ q),\\
  \norm{x}{X^{\Z}_{\theta,q}} &= \Phi^{\Z}_{\theta,q}\bigl(k \mapsto K(2^k,x)\bigr)
  && (\texttt{X.discreteKMethod}\ \theta\ q),
\end{align*}
and the finite \(J\)-norm (\texttt{jInfNorm X \(\theta\) q x}) is
\[
  \cJ_{\theta,q}(x; X) = \inf\Bigl\{\Phi^{\Z}_{\theta,q}\bigl(k \mapsto J(2^k,u_k)\bigr)
  : u\colon\Z\to\gamma \text{ finitely supported},\ \textstyle\sum_k u_k = x\Bigr\}.
\]
We call such \(u\) an \emph{admissible decomposition} of \(x\), and the value of
\(\Phi^{\Z}_{\theta,q}\) its \emph{cost}.

\paragraph{The data.}
Let \(\iota\) be a finite type. For \(i\in\iota\) let \(\alpha_i\) be a commutative group with a
decomposable modulus and \(A_i = (A_{i0}, A_{i1})\) a couple of solid quasinorms on
\(\alpha_i\) (\texttt{A i}). Let \(\beta\) be a commutative monoid with a modulus
\(\abs{\cdot}\colon\beta\to M\) and \(B = (B_0, B_1)\) a couple of quasinorms on \(\beta\).
For \(f \in \prod_i\alpha_i\), \(i\in\iota\) and \(x\in\alpha_i\), write \(\upd{f}{i}{x}\) for
\texttt{Function.update f i x}. The map \(T\colon\prod_i\alpha_i\to\beta\) is
\emph{multisubadditive} (\texttt{MultisubadditiveMap}):
\[
  \abs{T(\upd{f}{i}{x+y})} \le \abs{T(\upd{f}{i}{x})} + \abs{T(\upd{f}{i}{y})}.
\]
For finite sets \(s_i\subseteq\kappa_i\), \(\prod_i s_i\) denotes the finite set of
\(k\in\prod_i\kappa_i\) with \(k_i\in s_i\) for all \(i\) (\texttt{Fintype.piFinset s}).
Fix reals \(\alpha_0\) and \((\alpha_i)_{i\in\iota}\) (\texttt{c\(\alpha_0\)}, \texttt{c\(\alpha\)}),
and for \(\theta\in\R^\iota\) put \(\theta_0 = \alpha_0 + \sum_i \alpha_i\theta_i\).

\begin{definition}[\(\Omega\), corrected]\label{def:Omega}
  \(\theta\in\Omega\) if \(\theta_0 \ge 0\) and there are \(q_0, q_i \in (0,\infty]\) and
  \(C<\infty\) with
  \[
    \norm{Ta}{B_{\theta_0,q_0}} \le C\prod_i \norm{a_i}{(A_i)_{\theta_i,q_i}}
    \quad\text{whenever } a_i\in\Delta(A_i) \text{ and } \norm{a_i}{(A_i)_{\theta_i,q_i}} < \infty
    \text{ for all } i.
  \]
\end{definition}

The Lean definition \texttt{\(\Omega\)} differs only in that it allows \(q_0 = 0\) and
\(q_i = 0\); Remark~\ref{rem:q0} shows that this must be changed.

\begin{theorem}\label{thm:main}
  Assume moreover that the modulus of \(\beta\) is decomposable and that \(B_0, B_1\) are solid.
  Let \(\theta\in(0,1)^\iota\) with \(\theta_0\ge0\), and let \(C<\infty\) satisfy
  \begin{equation}\label{eq:ii}
    \norm{Ta}{B_{\theta_0,\infty}} \le C\prod_i \norm{a_i}{A_{i0}}^{1-\theta_i}\norm{a_i}{A_{i1}}^{\theta_i}
    \qquad\Bigl(a\in\prod_i\Delta(A_i)\Bigr). \tag{ii}
  \end{equation}
  Let \(C_B\) be the constant of \(B_{\theta_0,\infty}\) and \(p>0\) with \((2C_B)^p = 2\).
  Then there are constants \(D_i<\infty\), depending only on \(A_i\), \(\theta_i\) and \(p\), with
  \[
    \norm{Ta}{B_{\theta_0,\infty}} \le 2^{1/p}\, C \prod_i 2D_i\,\norm{a_i}{(A_i)_{\theta_i,p}}
    \qquad\Bigl(a\in\prod_i\Delta(A_i)\Bigr).
  \]
  In particular \(\theta\in\Omega\), with \(q_0 = \infty\) and \(q_i = p\).
\end{theorem}

\paragraph{Idea of the proof.}
This is Janson's argument (p.~295 of his paper). Decompose each \(a_i\) by a finite
\(J\)-decomposition, expand \(Ta\) by multisubadditivity into a sum over tuples of pieces,
bound each term by (ii) and by \(\norm{x}{0}^{1-\theta}\norm{x}{1}^{\theta} \le
2^{-k\theta}J(2^k,x)\), and sum the \(p\)-th powers (Aoki--Rolewicz). The sum over tuples of
a product factors into a product of \(\ell^p\)-sums. Janson's Lemma~1
(\texttt{jInfNorm\_le\_discreteKMethod}) then replaces \(\cJ\) by the discrete \(K\)-norm.

Two steps need more care than in Janson's two lines. First, \(M\) has no ordered addition and
\(B_{\theta_0,\infty}\) has constant \(C_B \ge 1\), so the expansion of \(Ta\) has to produce
a decomposition \(Ta = \sum_k w_k\) of the element itself (Lemma~\ref{lem:A2}); only then can
Aoki--Rolewicz be applied to a single finite sum. Second, (ii) is available only on
\(\prod_i\Delta(A_i)\), so the infimum over decompositions is taken over decompositions of
finite cost only (Lemma~\ref{lem:F}).

\section{Remarks on the statement}

\begin{remark}[\(q_0 = 0\)]\label{rem:q0}
  In the Lean definition of \(\Omega\) the exponent \(q_0\) may be \(0\). For \(q = 0\),
  \texttt{eLpNorm} returns \(0\) on a.e.\ strongly measurable functions
  (\texttt{eLpNorm\_exponent\_zero}). The integrand \(t\mapsto t^{-\theta}K(t,x)\) is
  measurable because \(t\mapsto K(t,x)\) is monotone (\texttt{monotone\_kNorm},
  \texttt{Monotone.measurable}). Hence \(\norm{x}{B_{\theta_0,0}} = 0\) for every \(x\), the
  boundedness condition holds with \(C = 0\), and the current \texttt{\(\Omega\)} equals
  \(\{\theta : \theta_0 \ge 0\}\). Then \texttt{mem\_\(\Omega\)\_of\_enorm\_le\_prod} is trivial
  and \texttt{convex\_\(\Omega\)} says nothing. The fix is to require \(q_0 \ne 0\), as in
  Janson (\(q\in(0,\infty]\)). The direction (i)\,\(\Rightarrow\)\,(ii) needs it as well,
  because it uses \(\norm{\cdot}{B_{\theta_0,\infty}} \lesssim \norm{\cdot}{B_{\theta_0,q_0}}\).
  Requiring \(q_i\ne0\) matches Janson. It is not needed for soundness, since \(q_i = 0\) only
  makes the bound harder to satisfy.
\end{remark}

\begin{remark}[ranges of \(\theta\)]
  Without \(\theta_0 \ge 0\) the theorem is false: \(T = 0\) satisfies (ii) for every
  \(\theta\). In Janson, condition (ii) contains \(\theta_0\in[0,1]\). The restriction
  \(\theta_i\in(0,1)\) is needed by \texttt{jInfNorm\_le\_discreteKMethod}. At
  \(\theta_i\in\{0,1\}\), Janson uses the convention \(A_{i,0,q} = A_{i0}\), which differs
  from the project's \texttt{kMethod}. Lemma~\ref{lem:C} and Proposition~\ref{prop:E} hold for
  \(\theta_i\in[0,1]\).
\end{remark}

\begin{remark}[solidity of \(B\)]
  Without solidity of \(B_0, B_1\), multisubadditivity in \(M\) says nothing about
  \(\norm{T\cdot}{}\), and the theorem fails.
\end{remark}

\begin{example}
  Let \(\iota\) be a singleton, \(\theta\in(0,1)\), \(\alpha_0 = 0\) and \(\alpha_1 = 1\), so
  that \(\theta_0 = \theta\).
  \begin{itemize}
    \item Let \(\alpha\) be the measurable real functions on \((0,1]\) modulo a.e.\ equality,
      with the pointwise modulus, \(A_0 = L^1\) and \(A_1 = L^\infty\) (extended norms). Both
      are solid.
    \item Let \(\beta = \R\) with \(B_0 = B_1 = \abs{\cdot}\) (absolute value), and the
      modulus \(\mu\colon\R\to[0,\infty]\), \(\mu(0) = 0\) and \(\mu(x) = \infty\) for
      \(x\ne0\). It is decomposable. Indeed, \(\mu(x+y)\le\mu(x)+\mu(y)\) because
      \(x+y\ne0\) forces \(x\ne0\) or \(y\ne0\). If \(\mu(x)\le\mu(y_0)+\mu(y_1)\) and
      \(x\ne0\), then some \(y_j\), say \(y_0\), is nonzero, and \(x = x + 0\) has the four
      required properties. If \(x = 0\), take \(0+0\). The absolute value is not solid for
      \(\mu\): \(\mu(1) = \mu(2)\) but \(1<2\).
    \item Let \(G(a) = \norm{a}{1}^{1-\theta}\norm{a}{\infty}^{\theta}\) and \(T(a) = G(a)\)
      if \(G(a)<\infty\), \(T(a) = 1\) otherwise. For \(a\ne0\) both norms are positive, so
      \(T(a)\ne0\); and \(T(0) = 0\). Hence \(\mu(T(a))\) is \(\infty\) for \(a\ne0\) and
      \(0\) for \(a = 0\). So \(T\) is subadditive for \(\mu\), because \(a+b\ne0\) forces
      \(a\ne0\) or \(b\ne0\).
  \end{itemize}
  For \(x\in\R\) we have \(K(t,x;B) = \min(1,t)\abs{x}\), so
  \(\norm{x}{B_{\theta,q_0}} = c_{q_0}\abs{x}\) with
  \(c_{q_0} = \Phi_{\theta,q_0}(\min(1,\cdot))\in(0,\infty)\) for \(q_0\in(0,\infty]\), and
  \(c_\infty = 1\). On \(\Delta = L^1\cap L^\infty\), \(G\) is finite, so (ii) holds with
  \(C = 1\). Now let \(f_\varepsilon = x^{\theta-1}\one_{[\varepsilon,1]}\). Then
  \(\norm{f_\varepsilon}{1} = (1-\varepsilon^\theta)/\theta\) and
  \(\norm{f_\varepsilon}{\infty} = \varepsilon^{\theta-1}\), so
  \(G(f_\varepsilon) \sim \theta^{\theta-1}\varepsilon^{-\theta(1-\theta)}\). On the other
  hand, splitting \(f_\varepsilon\) at height \(t^{\theta-1}\) gives
  \(K(t,f_\varepsilon) \le (1+1/\theta)t^\theta\). Together with
  \(K(t,f_\varepsilon)\le\norm{f_\varepsilon}{1}\le1/\theta\) and
  \(K(t,f_\varepsilon)\le t\varepsilon^{\theta-1}\), this gives
  \(\norm{f_\varepsilon}{\theta,q_1} = O(\log(1/\varepsilon)^{1/q_1})\). So
  \(\norm{Tf_\varepsilon}{B_{\theta,q_0}}/\norm{f_\varepsilon}{\theta,q_1}\to\infty\) for all
  \(q_0,q_1\in(0,\infty]\), and \(\theta\notin\Omega\).
\end{example}

\section{Auxiliary lemmas}

\begin{lemma}[Aoki--Rolewicz for finite sums]\label{lem:B}
  Let \(X\) be a quasinorm on a commutative monoid and \(p>0\) with \((2C_X)^p = 2\). For
  every finite family \((w_j)_{j\in s}\),
  \(\norm{\sum_{j\in s} w_j}{X}^p \le 2\sum_{j\in s}\norm{w_j}{X}^p\).
\end{lemma}

\begin{proof}
  This is \texttt{enorm\_sum\_rpow\_le\_two\_mul\_sum} in \texttt{Scratches/AokiRolewicz.lean}.
  It is inequality~(5) of Bergh--L\"ofstr\"om, Lemma~3.10.1, which rests on
  \(\norm{x+y}{X}^p \le 2\max(\norm{x}{X}^p,\norm{y}{X}^p)\).
\end{proof}

\begin{lemma}\label{lem:p}
  For \(1\le C<\infty\) there is \(p>0\) with \((2C)^p = 2\).
\end{lemma}

\begin{proof}
  Put \(r = 2C\in[2,\infty)\) and \(p = \log_r 2\). Then \(p>0\) because \(r>1\) and \(2>1\),
  and \(r^p = 2\) by the definition of the logarithm.
\end{proof}

\begin{lemma}\label{lem:C}
  Let \(X\) be a couple, \(\theta\in[0,1]\), \(k\in\Z\) and \(x\in\gamma\). Then
  \(\norm{x}{X_0}^{1-\theta}\norm{x}{X_1}^{\theta} \le 2^{-k\theta}J(2^k,x)\).
\end{lemma}

\begin{proof}
  Put \(t = 2^k\in(0,\infty)\). Since \(\norm{x}{X_0}\le J(t,x)\) and
  \(t\norm{x}{X_1}\le J(t,x)\),
  \[
    t^\theta\norm{x}{X_0}^{1-\theta}\norm{x}{X_1}^{\theta}
    = \norm{x}{X_0}^{1-\theta}\bigl(t\norm{x}{X_1}\bigr)^{\theta}
    \le J(t,x)^{1-\theta}J(t,x)^{\theta} = J(t,x).
  \]
  Multiplying by \(t^{-\theta} = 2^{-k\theta}\) gives the claim.
\end{proof}

\begin{lemma}\label{lem:D}
  Let \(0<p<\infty\), \(f\colon\Z\to[0,\infty]\) and \(s\subseteq\Z\) finite. Then
  \(\sum_{k\in s}(2^{-k\theta}f_k)^p \le \Phi^{\Z}_{\theta,p}(f)^p\).
\end{lemma}

\begin{proof}
  For \(0<p<\infty\), \(\Phi^{\Z}_{\theta,p}(f)^p = \sum_{k\in\Z}(2^{-k\theta}f_k)^p\). A
  finite partial sum of a series with terms in \([0,\infty]\) is at most the full sum.
\end{proof}

\begin{corollary}\label{cor:D}
  If \(0<p<\infty\) and \(\Phi^{\Z}_{\theta,p}(f)<\infty\), then \(f_k<\infty\) for every
  \(k\in\Z\).
\end{corollary}

\begin{proof}
  By Lemma~\ref{lem:D} with \(s = \{k\}\), \((2^{-k\theta}f_k)^p<\infty\). Since \(p>0\) and
  \(2^{-k\theta}\ne0\), \(f_k<\infty\).
\end{proof}

The next two lemmas expand \(Ta\). Since \(M\) has no ordered addition, one cannot first bound
\(\abs{T(\sum\cdots)}\) by a sum in \(M\) and decompose afterwards. Also,
\texttt{exists\_decomp} only splits along two moduli of elements of \(\beta\). So the element
\(z\) is decomposed along with the induction. The nonemptiness assumptions are needed because
nothing is known about \(T\) at \(0\).

\begin{lemma}[one coordinate]\label{lem:A1}
  Assume the modulus of \(\beta\) is decomposable. Let \(f\in\prod_i\alpha_i\),
  \(i\in\iota\), \(t\subseteq\kappa\) finite and nonempty, \(v\colon\kappa\to\alpha_i\), and
  \(z\in\beta\) with \(\abs{z}\le\abs{T(\upd{f}{i}{\sum_{a\in t}v_a})}\). Then
  \(z = \sum_{a\in t}z_a\) with \(\abs{z_a}\le\abs{T(\upd{f}{i}{v_a})}\) for \(a\in t\).
\end{lemma}

\begin{proof}
  By induction on the nonempty set \(t\). If \(t = \{a\}\), take \(z_a = z\). Otherwise
  \(t = \{a\}\cup t'\) with \(a\notin t'\) and \(t'\) nonempty. Then
  \(\sum_t v = v_a + \sum_{t'}v\), and multisubadditivity together with transitivity gives
  \[
    \abs{z} \le \abs{T(\upd{f}{i}{v_a})} + \abs{T(\upd{f}{i}{\textstyle\sum_{t'}v})}.
  \]
  By decomposability, \(z = z_a + z'\) with \(\abs{z_a}\le\abs{T(\upd{f}{i}{v_a})}\) and
  \(\abs{z'}\le\abs{T(\upd{f}{i}{\sum_{t'}v})}\). The induction hypothesis applied to
  \(z'\) and \(t'\) gives \(z' = \sum_{b\in t'}z_b\) with the required bounds, so
  \(z = \sum_{b\in t}z_b\).
\end{proof}

\begin{lemma}[all coordinates]\label{lem:A2}
  Assume the modulus of \(\beta\) is decomposable. Let \(s_i\subseteq\kappa_i\) be finite and
  nonempty, \(u_i\colon\kappa_i\to\alpha_i\), and \(z\in\beta\) with
  \(\abs{z}\le\abs{T(i\mapsto\sum_{k\in s_i}u_{i,k})}\). Then there is
  \(w\colon\prod_i\kappa_i\to\beta\) with
  \[
    z = \sum_{k\in\prod_i s_i} w_k
    \qquad\text{and}\qquad
    \abs{w_k}\le\abs{T(i\mapsto u_{i,k_i})}\quad\Bigl(k\in\prod_i s_i\Bigr).
  \]
\end{lemma}

\begin{proof}
  For \(S\subseteq\iota\) let \(P(S)\) be the statement: for every family \((t_i)\) of
  nonempty finite sets \(t_i\subseteq\kappa_i\) such that \(t_i\) is a singleton for
  \(i\notin S\), and every \(z\) with \(\abs{z}\le\abs{T(i\mapsto\sum_{t_i}u_i)}\), the
  conclusion holds with \(t\) in place of \(s\). The lemma is \(P(\iota)\). We prove \(P(S)\)
  by induction on \(S\).

  If \(S = \emptyset\), then \(t_i = \{c_i\}\) for all \(i\), so \(\prod_i t_i = \{c\}\), the
  hypothesis reads \(\abs{z}\le\abs{T(i\mapsto u_{i,c_i})}\), and \(w_c = z\) works.

  Let \(j\notin S\), assume \(P(S)\), and let \(t\) be singletons outside \(S\cup\{j\}\). Put
  \(f = (i\mapsto\sum_{t_i}u_i)\), so that \(f = \upd{f}{j}{\sum_{a\in t_j}u_{j,a}}\).
  Lemma~\ref{lem:A1} at the coordinate \(j\) gives \(z = \sum_{a\in t_j}z_a\) with
  \(\abs{z_a}\le\abs{T(\upd{f}{j}{u_{j,a}})}\). For \(a\in t_j\) let
  \(t^a = \upd{t}{j}{\{a\}}\). Then \(\upd{f}{j}{u_{j,a}} = (i\mapsto\sum_{t^a_i}u_i)\), and
  \(t^a_i\) is a singleton for \(i\notin S\). So \(P(S)\) gives \(w^a\) with
  \(z_a = \sum_{k\in\prod t^a}w^a_k\) and the required bounds. Now
  \(\prod_i t^a_i = \{k\in\prod_i t_i : k_j = a\}\), and these sets partition \(\prod_i t_i\)
  as \(a\) ranges over \(t_j\). Put \(w_k = w^{k_j}_k\). Summing over the fibres of
  \(k\mapsto k_j\),
  \[
    \sum_{k\in\prod t}w_k = \sum_{a\in t_j}\ \sum_{k\in\prod t,\ k_j = a}w^a_k
    = \sum_{a\in t_j}z_a = z,
  \]
  and for \(k\in\prod t\) we have \(k\in\prod t^{k_j}\), hence
  \(\abs{w_k}\le\abs{T(i\mapsto u_{i,k_i})}\).
\end{proof}

\begin{lemma}\label{lem:F}
  Let \(f_i\colon\kappa_i\to[0,\infty]\) for \(i\in\iota\), each taking at least one finite
  value. Let \(K<\infty\) and \(a\in[0,\infty]\) satisfy \(a\le K\prod_i f_i(g_i)\) for every
  \(g\in\prod_i\kappa_i\) with \(f_i(g_i)<\infty\) for all \(i\). Then
  \(a\le K\prod_i\inf_j f_i(j)\).
\end{lemma}

\begin{proof}
  Let \(m_i = \inf_j f_i(j)\). Infinite values do not lower the infimum and a finite value
  exists, so \(m_i\) is the infimum of the finite values of \(f_i\), and \(m_i<\infty\).
  Choose \(g^n_i\) with \(f_i(g^n_i)<\infty\) and \(f_i(g^n_i)\to m_i\). Multiplication on
  \([0,\infty]\) is continuous at pairs of finite numbers, and all \(m_i\) and \(K\) are
  finite, so \(K\prod_i f_i(g^n_i)\to K\prod_i m_i\). Since \(a\le K\prod_i f_i(g^n_i)\) for
  every \(n\), the claim follows.
\end{proof}

The finiteness of the limits matters. An algebraic induction based on
\(a\le(\inf_i f_i)(\inf_j g_j)\) runs into \(\infty\cdot0 = 0\) at tuples where one factor
is infinite, and those are exactly the tuples where (ii) is not available.

\begin{lemma}\label{lem:solid}
  Let \(X_0, X_1\) be solid quasinorms on \(\gamma\) for a decomposable modulus. Then
  \(X_{\theta,q}\) is solid for all \(\theta, q\).
\end{lemma}

\begin{proof}
  Let \(\abs{x}\le\abs{y}\) and \(t\in[0,\infty]\). For every decomposition
  \(y = y_0 + y_1\), \(\abs{x}\le\abs{y}\le\abs{y_0}+\abs{y_1}\), so decomposability gives
  \(x = x_0 + x_1\) with \(\abs{x_j}\le\abs{y_j}\). By solidity,
  \(K(t,x)\le\norm{x_0}{X_0} + t\norm{x_1}{X_1}\le\norm{y_0}{X_0} + t\norm{y_1}{X_1}\), and
  taking the infimum gives \(K(t,x)\le K(t,y)\). The function \(t\mapsto K(t,x)\) is monotone,
  hence measurable, and \(\Phi_{\theta,q}\) is monotone on measurable functions. Hence
  \(\norm{x}{X_{\theta,q}}\le\norm{y}{X_{\theta,q}}\).
\end{proof}

\section{The key estimate}

\begin{proposition}\label{prop:E}
  Assume the modulus of \(\beta\) is decomposable and \(B_0, B_1\) are solid. Let
  \(\theta\in[0,1]^\iota\), assume (ii) with constant \(C<\infty\), and let \(p>0\) with
  \((2C_B)^p = 2\). Then for all \(a\in\prod_i\Delta(A_i)\),
  \[
    \norm{Ta}{B_{\theta_0,\infty}} \le 2^{1/p}\,C\prod_i\cJ_{\theta_i,p}(a_i;A_i).
  \]
\end{proposition}

\begin{proof}
  \emph{Step 1: one choice of decompositions.} For each \(i\) let \(u_i\colon\Z\to\alpha_i\) be
  an admissible decomposition of \(a_i\) with finite cost
  \(c_i = \Phi^{\Z}_{\theta_i,p}(k\mapsto J(2^k,u_{i,k}))\). We show
  \begin{equation}\label{eq:fixed-u}
    \norm{Ta}{B_{\theta_0,\infty}} \le 2^{1/p}\,C\prod_i c_i .
  \end{equation}
  By Corollary~\ref{cor:D}, \(J(2^k,u_{i,k})<\infty\) for every \(k\). Since
  \(\norm{u_{i,k}}{A_{i0}}\le J(2^k,u_{i,k})\) and
  \(2^k\norm{u_{i,k}}{A_{i1}}\le J(2^k,u_{i,k})\) with \(2^k>0\), every \(u_{i,k}\) lies in
  \(\Delta(A_i)\).

  Let \(s_i\) be the support of \(u_i\) together with \(0\). It is finite and nonempty, and
  \(\sum_{k\in s_i}u_{i,k} = a_i\). Lemma~\ref{lem:A2} with \(z = Ta\) gives
  \(Ta = \sum_{k\in\prod_i s_i}w_k\) with \(\abs{w_k}\le\abs{Tu_k}\), where
  \(u_k = (u_{i,k_i})_i\in\prod_i\Delta(A_i)\). With all sums over \(k\in\prod_i s_i\),
  \begin{align*}
    \norm{Ta}{B_{\theta_0,\infty}}^p
    &\le 2\sum_k\norm{w_k}{B_{\theta_0,\infty}}^p
    && \text{(Lemma~\ref{lem:B})}\\
    &\le 2\sum_k\norm{Tu_k}{B_{\theta_0,\infty}}^p
    && \text{(Lemma~\ref{lem:solid})}\\
    &\le 2C^p\sum_k\prod_i\bigl(\norm{u_{i,k_i}}{A_{i0}}^{1-\theta_i}\norm{u_{i,k_i}}{A_{i1}}^{\theta_i}\bigr)^p
    && \text{((ii), as } u_k\in\textstyle\prod_i\Delta(A_i))\\
    &\le 2C^p\sum_k\prod_i\bigl(2^{-k_i\theta_i}J(2^{k_i},u_{i,k_i})\bigr)^p
    && \text{(Lemma~\ref{lem:C})}\\
    &= 2C^p\prod_i\sum_{k\in s_i}\bigl(2^{-k\theta_i}J(2^{k},u_{i,k})\bigr)^p
    && \text{(distributivity)}\\
    &\le 2C^p\prod_i c_i^p
    && \text{(Lemma~\ref{lem:D}).}
  \end{align*}
  Taking \(p\)-th roots gives \eqref{eq:fixed-u}.

  \emph{Step 2: the infimum.} For each \(i\), let \(f_i(u)\) be the cost of \(u\) if \(u\) is
  an admissible decomposition of \(a_i\), and \(\infty\) otherwise; then
  \(\inf_u f_i(u) = \cJ_{\theta_i,p}(a_i;A_i)\). The decomposition with \(u_0 = a_i\) and
  \(u_k = 0\) for \(k\ne0\) has cost \(J(1,a_i)\), which is finite since
  \(a_i\in\Delta(A_i)\). By Step~1, the hypothesis of Lemma~\ref{lem:F} holds with
  \(a = \norm{Ta}{B_{\theta_0,\infty}}\) and \(K = 2^{1/p}C<\infty\), and the lemma gives the
  claim.
\end{proof}

\section{Proof of Theorem~\ref{thm:main}}

Since \(1\le C_B<\infty\), Lemma~\ref{lem:p} provides \(p\). Let
\(a\in\prod_i\Delta(A_i)\). Proposition~\ref{prop:E} gives
\(\norm{Ta}{B_{\theta_0,\infty}}\le2^{1/p}C\prod_i\cJ_{\theta_i,p}(a_i;A_i)\). Since
\(\theta_i\in(0,1)\), \(a_i\in\Delta(A_i)\) and \(A_{i0}, A_{i1}\) are solid,
\texttt{jInfNorm\_le\_discreteKMethod} gives
\[
  \cJ_{\theta_i,p}(a_i;A_i)\le2^{\min(\theta_i,1-\theta_i)}\norm{a_i}{(A_i)^{\Z}_{\theta_i,p}}
  \le2\norm{a_i}{(A_i)^{\Z}_{\theta_i,p}}.
\]
By \texttt{discreteKMethod\_equiv\_kmethod} there is \(D_i<\infty\) with
\(\norm{x}{(A_i)^{\Z}_{\theta_i,p}}\le D_i\norm{x}{(A_i)_{\theta_i,p}}\) for all \(x\). This is the
half \texttt{.ge} of the equivalence, since \(X\le Y\) means
\(\exists D<\infty,\ \forall x,\ \norm{x}{Y}\le D\norm{x}{X}\). Combining the three bounds
proves the inequality of the theorem. Its constant is finite, \(q_0 = \infty\) and
\(q_i = p\) are nonzero, and \(\theta_0\ge0\). So \(\theta\in\Omega\). \qed

\medskip
The bound holds on all of \(\prod_i\Delta(A_i)\), without
\(\norm{a_i}{(A_i)_{\theta_i,p}}<\infty\). The hypothesis \(\alpha_i\ne0\)
(\texttt{h\(\alpha\)}) is not used.

\section{Lean plan}
\raggedright

\paragraph{Statement changes.}
\begin{itemize}
  \item \texttt{\(\Omega\)}: add \texttt{q\(_0\) \(\ne\) 0 \(\wedge\) \(\forall\) i, q i \(\ne\) 0}.
  \item \texttt{mem\_\(\Omega\)\_of\_enorm\_le\_prod} (and \texttt{mem\_\(\Omega\)\_iff}): add
    \texttt{(h\(\theta\) : \(\forall\) i, \(\theta\) i \(\in\) Ioo 0 1)},
    \texttt{(h\(\theta_0\) : 0 \(\le\) c\(\alpha_0\) + \(\sum\) i, c\(\alpha\) i * \(\theta\) i)} and
    \texttt{[Modulus.IsDecomposable \(\beta\) M] [B.fst.IsSolid] [B.snd.IsSolid]}.
\end{itemize}

\paragraph{New declarations} (suggested names).
\begin{itemize}
  \item Lemma~\ref{lem:B}: move the \texttt{Dyadic} section and
    \texttt{enorm\_add\_rpow\_le\_two\_mul\_max}, \texttt{enorm\_sum\_rpow\_le\_of\_weights},
    \texttt{enorm\_sum\_rpow\_le\_two\_mul\_sum} from \texttt{Scratches/AokiRolewicz.lean} (not
    built, not a \texttt{module}) into \texttt{MultilinearInterpolation/AokiRolewicz.lean}.
  \item Lemma~\ref{lem:p}: \texttt{EQuasinorm.exists\_rpow\_two\_mul\_C\_eq\_two}, next to
    \texttt{pos\_of\_rpow\_two\_mul\_C\_eq\_two}. Uses \texttt{Real.logb\_pos},
    \texttt{Real.rpow\_logb}, \texttt{ENNReal.ofReal\_rpow\_of\_pos},
    \texttt{ENNReal.ofReal\_toReal}. (The \texttt{C\_eq\_rpow\_iff} cited in the scratch
    docstring does not exist.)
  \item Lemma~\ref{lem:C}: no new lemma. Use
    \texttt{enorm\_rpow\_mul\_enorm\_rpow\_le\_rpow\_neg\_mul\_jNorm} with \(t = 2^k\) and
    \texttt{ENNReal.rpow\_intCast\_mul}.
  \item Lemma~\ref{lem:D}: general form
    \(\sum_{k\in s}\norm{g_k}{}^{p} \le (\texttt{eLpNorm}\ g\ p\ \texttt{count})^{p}\) in the
    project's \texttt{Mathlib/\dots/LpSeminorm/Count.lean}, via
    \texttt{eLpNorm\_eq\_lintegral\_rpow\_enorm\_toReal}, \texttt{ENNReal.rpow\_inv\_rpow},
    \texttt{lintegral\_count}, \texttt{ENNReal.sum\_le\_tsum}.
  \item Lemma~\ref{lem:A1}: \texttt{MultisubadditiveMap.exists\_sum\_eq\_of\_modulus\_le\_update};
    \texttt{Finset.Nonempty.cons\_induction}, \texttt{Finset.sum\_cons},
    \texttt{Modulus.IsDecomposable.exists\_decomp}.
  \item Lemma~\ref{lem:A2}: \texttt{MultisubadditiveMap.exists\_sum\_eq\_of\_modulus\_le};
    \texttt{Finset.induction\_on}, \texttt{Fintype.piFinset\_singleton},
    \texttt{Fintype.piFinset\_update\_singleton\_eq\_filter\_piFinset\_eq},
    \texttt{Finset.sum\_fiberwise\_of\_maps\_to}. Needs
    \texttt{[\(\forall\) i, DecidableEq (\(\kappa\) i)]}. The hard step is identifying
    \(\upd{f}{j}{u_{j,a}}\) with the sum over \(t^a\) (dependent \texttt{Function.update}).
  \item Lemma~\ref{lem:F}: e.g.\ \texttt{ENNReal.le\_mul\_prod\_iInf}; via
    \texttt{iInf\_ne\_top\_subtype}, \texttt{exists\_seq\_tendsto\_sInf},
    \texttt{ENNReal.tendsto\_finsetProd\_of\_ne\_top}, \texttt{ge\_of\_tendsto'}.
  \item Lemma~\ref{lem:solid}: \texttt{phiFunctional\_mono} (via \texttt{eLpNorm\_mono\_enorm},
    measurability from \texttt{Monotone.measurable}) and an instance
    \texttt{IsSolid.kMethod} (via \texttt{kNorm\_le\_kNorm\_of\_modulus\_le},
    \texttt{monotone\_kNorm}).
  \item Proposition~\ref{prop:E}: e.g.\ \texttt{enorm\_le\_mul\_prod\_jInfNorm}; also uses
    \texttt{mem\_finiteLocus\_inf}, \texttt{finsum\_eq\_sum\_of\_support\_subset},
    \texttt{Finset.prod\_univ\_sum}, \texttt{ENNReal.prod\_rpow\_of\_nonneg},
    \texttt{discretePhiFunctional\_single}.
\end{itemize}

\end{document}
